% 7.2：沿用既有生产解；全域控制位置核验及事件定位与code/p23.py一致。
\subsection{模型求解}
\label{sec:p3-solution}

\subsubsection{时间推进与临界时刻定位}

前$3\ \mathrm{h}$沿用第\ref{sec:p2-solution}节的数值设置，
此后取$\Delta t_n=1\ \mathrm{s}$，网格、双场迭代与时间格式不变，
全程仍仅最初两步作隐式启动。

每步收敛后，在全部计算节点上求
$C_{\max}^{n}=\max_{0\le i\le N}C_i^{n}$，
并与中心含水率比较。
全过程核验表明，二者之差始终不超过$10^{-12}\cunit$，
故本次计算中可用中心含水率定位全域达标时刻。
当首次出现$C_{\max}^{n}\ge C_*>C_{\max}^{n+1}$时，认为在该区间内水分浓度持续减小，
临界时刻位于$[t^n,t^{n+1})$内。
对该步的未舍入结果作线性插值，得到
\begin{equation}
\begin{gathered}
  \lambda=\frac{C_{\max}^{n}-C_*}
                 {C_{\max}^{n}-C_{\max}^{n+1}},
  \qquad t_*\approx t^n+\lambda\Delta t_n,\\
  \boldsymbol C(t_*)\approx
    (1-\lambda)\boldsymbol C^n+\lambda\boldsymbol C^{n+1}.
\end{gathered}
\label{eq:p3-event-interpolation}
\end{equation}
该式同时给出临界时刻及径向含水率分布，规定的$60\ \mathrm{s}$间隔仅用于结果采样。

\subsubsection{烘干时长与计算结果}

计算得到，在前述环境延拓条件下，临界烘干时间为
\begin{equation*}
  \boxed{t_*\approx57.4721\ \mathrm{h}}.
\end{equation*}
每隔$6\ \mathrm{h}$提取$r=0,0.5,1,1.5,2\ \mathrm{cm}$处的含水率，
并由式\eqref{eq:p3-event-interpolation}补充临界时刻的分布，结果见表\ref{tab:p3C}。

\begin{table}[!htbp]
  \centering
  \caption{药材烘干过程的水分浓度（单位：$\mathrm{kg/kg}$）}
  \label{tab:p3C}
  \begin{tabularx}{\textwidth}{@{}C{0.22\textwidth}YYYYY@{}}
    \toprule
    \multirow{2}{*}{时间/$\mathrm{h}$}&\multicolumn{5}{c}{到药材中心的距离/$\mathrm{cm}$}\\
    \cmidrule(lr){2-6}
    &0&0.5&1&1.5&2\\
    \midrule
    6  &1.0170&0.9881&0.9015&0.7550&0.5328\\
    12 &0.4561&0.4432&0.4031&0.3297&0.1642\\
    18 &0.2988&0.2915&0.2687&0.2253&0.0845\\
    24 &0.2378&0.2327&0.2167&0.1855&0.0665\\
    30 &0.2058&0.2018&0.1892&0.1641&0.0598\\
    36 &0.1859&0.1825&0.1718&0.1504&0.0566\\
    42 &0.1721&0.1691&0.1597&0.1407&0.0549\\
    48 &0.1618&0.1592&0.1507&0.1334&0.0537\\
    54 &0.1539&0.1514&0.1436&0.1277&0.0530\\
    \mbox{\textbf{57.4721（结束）}}
       &\textbf{0.1500}&\textbf{0.1477}&\textbf{0.1402}&\textbf{0.1249}&\textbf{0.0527}\\
    \bottomrule
  \end{tabularx}
\end{table}

表末行为$C_{\max}=C_*$的临界状态；继续推进至其后的首个整分钟\textbf{$206940\ \mathrm{s}$}，
此时未舍入的$C_{\max}<C_*$。
全程每隔$60\ \mathrm{s}$、径向每隔$0.1\ \mathrm{cm}$采样，以四位小数写入\texttt{result3.xlsx}。
